% GENERATED FILE. Edit canonical lessons, then run npm run generate:books.
% canonical-source-hash: fnv1a64:9ec024df478a8680
% canonical-lessons: PA-R170-date-selector-again

\chapter{Return to the Fictional Date Selector}
\label{ch:pa-date-selector-r4}
\hlchaptermodality{170}

\begin{chapteropening}
\textbf{By the end of this chapter:} \emph{I can retrieve the already-taught two-card date selector after a long gap, with the fictional bank visible and without claiming scored A1 writing.}

Recall that \pa{ਕ} selects fictional date A and \pa{ਖ} selects fictional date B without adding new Punjabi.
\end{chapteropening}

\section[\texorpdfstring{\pa{ਕ} · \pa{ਖ} (the two familiar …)}{ਕ · ਖ (the two familiar …)}]{The old two-card selector}

\label{lesson:PA-R170-date-selector-again}



\begin{quote}
Look at the two old cards: A \textbf{\pa{੧੫}/\pa{੦੧}/\pa{੨੦੨੫}} and B \textbf{\pa{੨੫}/\pa{੦੨}/\pa{੨੦੨੫}}. The mark \textbf{\pa{ਕ}} asks for A; \textbf{\pa{ਖ}} asks for B. These marks are selectors, not month names.
\end{quote}

\subsection*{Your turn}
Keep both cards and the selector map visible. Point to \textbf{\pa{ਖ}}, then to B. Write the letter B beside \textbf{\pa{ਖ}}. Do the same for \textbf{\pa{ਕ}} and A. Compare each link with the printed map before trying the pair once more with the map covered. If one link differs, uncover the map and repair that link. No new date is needed.

\begin{tcolorbox}[breakable,colback=teal!4,colframe=teal!35!black,title={Before you move on}]
Cover the map for a few seconds. Point to the card requested by \textbf{\pa{ਕ}}, then the card requested by \textbf{\pa{ਖ}}. Uncover the map and check A, then B. The bank remains available for a gentle repair; this is not independent scored writing.
\end{tcolorbox}
